% !TEX root = ../../cookbook.tex
\newpage
\recipe{Ciabatta bread}{\href{https://www.youtube.com/watch?v=jA95tUK2RZU&t=4s}{Brian Lagerstrom's recipe}, Makes 2 large ciabattas.}{}{}

\begin{multicols}{2}
	\subsection*{Ingredients}
	
	\subsubsection*{For the Biga}
	\begin{ingredients}
		\item 175 g warm water
		\item \sfrac{1}{4} tsp active dry yeast
		\item 225 g bread flour (11.7\% protein)
	\end{ingredients}
	
	\subsubsection*{For the Autolyse}
	\begin{ingredients}
		\item 180 g warm water
		\item 250 g bread flour
	\end{ingredients}
	\columnbreak
	\subsection*{ }
	\subsubsection*{For the Final Mix}
	\begin{ingredients}
		\item 40 g warm water
		\item 5 g active dry yeast
		\item 10 g salt
		\item All of the biga
	\end{ingredients}
\end{multicols}

% --- Preparation ---
\textbf{Prepare the biga, ideally the day before:} \\
Combine the \textbf{water}, \textbf{yeast}, and \textbf{flour} in a tall container, mix well, and cover.
Leave to ferment at room temperature for 6--24 hours.

\textbf{Autolyse:} \\
Combine the \textbf{water} and \textbf{flour} in the bowl of a stand mixer.
Mix with the dough hook on low speed until fully combined, then cover with a clean kitchen towel and let rest for 30 minutes.

\textbf{Final mix:} \\
Add the \textbf{remaining ingredients} and the \textbf{biga} to the autolysed dough. Mix on low speed for 3 minutes until well combined, then on high speed for 5 minutes.
The dough should clean the sides of the bowl and become smooth, glossy, and elastic.

\textbf{Bulk fermentation:} \\
Transfer the dough to a large, well-oiled container, cover, and let rise for 30 minutes.

Perform a strength-building fold\footnote{strength-building fold: see the recipe \nameref{pizza_napolitana} (p. \pageref{pizza_napolitana}), footnote~\ref{fn:strength_building_fold}.} and let rise for another 30 minutes.

Perform another strength-building fold or, even better, \textit{laminate} the dough\footnote{Lamination is a technique used to strengthen dough. In pastry, it typically involves folding layers of butter into the dough to create many distinct layers, as in croissants. For bread dough, it consists of stretching the dough into a thin sheet on a wet work surface, folding the sides inward like a letter, and then rolling it back onto itself to form a tight dough ball. Brian Lagerstrom demonstrates this technique in the recipe video linked at the top of this page, below the title.}, then let rest for 1 hour.

\textbf{Shape the ciabattas:} \\
Flour the dough and the work surface, then gently turn the dough out, taking care to release it cleanly from the container without deflating it.
Divide the dough in half and, \textbf{very gently}, without pressing or further working it, shape each piece into a rectangular ciabatta loaf.
Transfer them to floured parchment paper, cover, and let rest for 30 minutes.

\textbf{Baking:} \\
Preheat the oven to 260°C.
Place the ciabattas in the oven together with a container of boiling water\footnote{The container of boiling water replaces the humidity control found in professional bread ovens. During the first stage of baking, the steam keeps the surface of the dough moist, delaying crust formation and allowing the loaves to expand more fully. Once the crust begins to form, the steam contributes to a thinner, crispier crust and promotes deeper browning.}, and bake for approximately 25 minutes, until deeply golden brown.

% --- Description ---
\begin{recipeblock}
	
\end{recipeblock}